\section{Experimental Design: Five Domains, 42{,}184 Swap Runs}
\label{sec:expdesign}

\subsection{Domains and seed prompts}
\label{sec:expdesign:domains}
We evaluate on five domains spanning single-hop factual recall in different relational structures and answer geometries (\cref{tab:domains}). The first four are the \emph{strong} domains used for cross-domain inference; \emph{Sounds} (6 entities, 12/30 pairs share answers, $K{=}6$) is reported separately due to known structural confounds (\cref{sec:results:negative}, \cref{sec:limitations}).

\begin{table*}[t]
\centering
\small
\caption{\textbf{Four evaluation domains.} Each row gives the seed-prompt template, the three semantic fields (input / intermediate / answer), the answer-set size $K$, the number of entities and non-identity swap pairs, and the total run count across all conditions (labeled, matched-random $\times 3$, 7 field-additivity variants, adaptive-$M$ rescue). Full per-condition counts in \cref{appx:specificity,appx:fieldadd,appx:msearch}.}
\label{tab:domains}
\begin{tabular}{llp{3.4cm}rrrr}
\toprule
Domain & Seed prompt template & Fields (in / int / ans) & $K$ & Entities & Pairs & Total runs \\
\midrule
USA       & ``capital of state containing $X$'' & city / state / capital       & 50 & 50 & 2{,}450 & $\approx$28{,}400 \\
Books     & ``character $X$ from book by''      & character / book / author    & 25 & 16 & 240     & $\approx$2{,}800  \\
Products  & ``product $X$ from company by''     & product / company / founder  & 25 & 12 & 132     & $\approx$1{,}700  \\
Paintings & ``painting $X$ by painter''          & painting / painter / first\_name & 20 & 10 & 90  & $\approx$1{,}200  \\
\midrule
\multicolumn{6}{l}{\textbf{Total}} & \textbf{41{,}784} \\
\bottomrule
\end{tabular}
\end{table*}


\subsection{Experimental matrix}
\label{sec:expdesign:matrix}
Each swap pair is run under: (a) labeled (3-field) baseline at $M{=}20$; (b) matched random control (3 replicates per pair); (c) field-additivity (7 variants per pair); (d) adaptive $M$-search on labeled and on the best field-add variant per missed pair. The total is \textbf{42{,}184 swap runs} across all conditions, domains, and $M$ values (per-condition counts in \cref{tab:domains}).

\subsection{Pre-registered evaluation choices}
\label{sec:expdesign:prereg}
Before analysis we fixed: (i) contrast group $\mathcal{A}$ = all other dataset answers; (ii) first-subword resolution for multi-token answers; (iii) explicit \texttt{answer\_field} decoupling intervention from measurement; (iv) generation parameters (T$=$0.3, $n{=}10$, frequency-penalty $2.0$, seed $42$); (v) sha256-seeded matched-control replicates; (vi) pair-level definition of Hit (any of 3 random replicates is a hit counts as a random-hit for that pair). Ranges, thresholds, and seeds were frozen prior to running the 42{,}184-run sweep.
